\documentclass{article}

\usepackage[utf8]{inputenc}
\usepackage[T1]{fontenc}
\usepackage[french]{babel}
\usepackage{url}
\usepackage{graphicx}
\usepackage{hyperref}
\usepackage{booktabs}
\usepackage{array}
\usepackage{float}
\usepackage{listings}
\usepackage{xcolor}
\usepackage{geometry}

\geometry{margin=2.5cm}

\setlength{\parskip}{0.8em}
\setlength{\parindent}{0pt}

% Style pour le code source
\lstset{
  basicstyle=\ttfamily\small,
  breaklines=true,
  frame=single,
  backgroundcolor=\color{gray!8},
  keywordstyle=\color{blue!70!black}\bfseries,
  commentstyle=\color{green!50!black}\itshape,
  stringstyle=\color{orange!80!black},
  numberstyle=\tiny\color{gray},
  numbers=left,
  numbersep=8pt,
  showstringspaces=false,
  tabsize=2,
}

% Numérotation des sections en chiffres romains
\renewcommand{\thesection}{\Roman{section}}

\title{\textbf{Annexes}\\[0.4em]
\large Développement d'une application web pour la gestion d'une école de musique}
\author{Baris Ozcelik}
\date{Année académique 2025--2026}

\begin{document}

\maketitle
\tableofcontents
\newpage

% ─────────────────────────────────────────────
\section{Schéma de la base de données (Prisma)}
% ─────────────────────────────────────────────

Le schéma suivant correspond au fichier \texttt{prisma/schema.prisma} du projet.
Il définit l'ensemble des modèles, relations et énumérations utilisés par l'application.

\begin{lstlisting}[language=,caption={prisma/schema.prisma}]
generator client {
  provider = "prisma-client"
  output   = "../generated/prisma"
}

datasource db {
  provider = "postgresql"
}

enum Role {
  ADMIN
  STUDENT
}

model User {
  id           String   @id @default(uuid())
  email        String   @unique
  passwordHash String?
  role         Role
  createdAt    DateTime @default(now())

  invitationToken   String?   @unique
  tokenExpiresAt    DateTime?
  isActive          Boolean   @default(false)

  resetToken          String?   @unique
  resetTokenExpiresAt DateTime?

  student Student?
}

enum PaymentMode {
  PER_SESSION
  MONTHLY
}

model Student {
  id          String      @id @default(uuid())
  userId      String      @unique
  firstName   String
  lastName    String
  phoneNumber String?
  balance            Int         @default(0)
  paymentMode        PaymentMode @default(PER_SESSION)
  pendingPaymentMode PaymentMode?

  user        User           @relation(fields: [userId], references: [id], onDelete: Cascade)
  enrollments Enrollment[]
  attendances Attendance[]
  accesses    ResourceAccess[]
  paymentRequested Boolean @default(false)
}

model Course {
  id        String   @id @default(uuid())
  title     String
  capacity  Int
  createdAt DateTime @default(now())

  enrollments         Enrollment[]
  sessions            Session[]
  inscriptionRequests InscriptionRequestCourse[]
}

model Enrollment {
  id        String   @id @default(uuid())
  studentId String
  courseId  String
  createdAt DateTime @default(now())

  student Student @relation(fields: [studentId], references: [id], onDelete: Cascade)
  course  Course  @relation(fields: [courseId], references: [id], onDelete: Cascade)
  @@unique([studentId, courseId])
}

model Event {
  id          String    @id @default(uuid())
  title       String
  description String?
  location    String
  startAt     DateTime
  endAt       DateTime?
  price       Int
  capacity    Int
  createdAt   DateTime  @default(now())

  tickets   Ticket[]
  order     Order[]
  scanToken ScanToken?
}

model Order {
  id              String   @id @default(uuid())
  email           String
  amount          Int
  status          String
  stripeSessionId String?  @unique
  eventId         String
  createdAt       DateTime @default(now())

  event   Event    @relation(fields: [eventId], references: [id], onDelete: Cascade)
  tickets Ticket[]
}

model Ticket {
  id      String    @id @default(uuid())
  qrCode  String    @unique
  usedAt  DateTime?
  orderId String
  eventId String
  createdAt DateTime @default(now())

  order Order @relation(fields: [orderId], references: [id], onDelete: Cascade)
  event Event @relation(fields: [eventId], references: [id], onDelete: Cascade)
}

enum SessionStatus {
  PLANNED
  CANCELLED
  COMPLETED
}

model Session {
  id        String        @id @default(uuid())
  courseId  String
  startsAt  DateTime
  endsAt    DateTime
  hours     Int
  status    SessionStatus @default(PLANNED)
  createdAt DateTime      @default(now())
  updatedAt DateTime      @updatedAt

  course      Course       @relation(fields: [courseId], references: [id], onDelete: Cascade)
  attendances Attendance[]
  @@index([courseId, startsAt])
}

enum AttendanceStatus {
  PRESENT
  ABSENT
  LATE
  EXCUSED
}

model Attendance {
  id        String           @id @default(uuid())
  sessionId String
  studentId String
  status    AttendanceStatus
  createdAt DateTime         @default(now())
  updatedAt DateTime         @updatedAt

  session Session @relation(fields: [sessionId], references: [id], onDelete: Cascade)
  student Student @relation(fields: [studentId], references: [id], onDelete: Cascade)
  @@unique([sessionId, studentId])
}

enum InscriptionStatus {
  PENDING
  ACCEPTED
  REJECTED
}

model InscriptionRequest {
  id              String            @id @default(uuid())
  firstName       String
  lastName        String
  email           String            @unique
  phoneNumber     String?
  isParent        Boolean           @default(false)
  parentFirstName String?
  parentLastName  String?
  message         String?
  status          InscriptionStatus @default(PENDING)
  createdAt       DateTime          @default(now())
  updatedAt       DateTime          @updatedAt

  courses InscriptionRequestCourse[]
}

model InscriptionRequestCourse {
  id                   String @id @default(uuid())
  inscriptionRequestId String
  courseId             String

  inscriptionRequest InscriptionRequest @relation(
    fields: [inscriptionRequestId], references: [id], onDelete: Cascade)
  course Course @relation(
    fields: [courseId], references: [id], onDelete: Cascade)
  @@unique([inscriptionRequestId, courseId])
}

model PricingConfig {
  id              String    @id @default(uuid())
  perSessionCents Int
  monthlyCents    Int
  effectiveFrom   DateTime
  appliedAt       DateTime?
  createdAt       DateTime  @default(now())
}

model ScanToken {
  id        String   @id @default(uuid())
  token     String   @unique @default(uuid())
  eventId   String   @unique
  expiresAt DateTime
  createdAt DateTime @default(now())

  event Event @relation(fields: [eventId], references: [id], onDelete: Cascade)
}

model Resource {
  id          String   @id @default(uuid())
  title       String
  description String?
  fileUrl     String
  fileName    String
  createdAt   DateTime @default(now())

  accesses ResourceAccess[]
}

model ResourceAccess {
  id         String   @id @default(uuid())
  resourceId String
  studentId  String
  createdAt  DateTime @default(now())

  resource Resource @relation(fields: [resourceId], references: [id], onDelete: Cascade)
  student  Student  @relation(fields: [studentId], references: [id], onDelete: Cascade)
  @@unique([resourceId, studentId])
}
\end{lstlisting}

\newpage

% ─────────────────────────────────────────────
\section{Variables d'environnement}
% ─────────────────────────────────────────────

Le tableau suivant liste l'ensemble des variables d'environnement requises pour le
bon fonctionnement de l'application, ainsi que leur rôle. Ces variables sont à
configurer dans un fichier \texttt{.env} en environnement local et dans le tableau
de bord \textit{Vercel} en production.

\begin{table}[H]
\centering
\renewcommand{\arraystretch}{1.4}
\begin{tabular}{>{\ttfamily}p{5.5cm} p{8.5cm}}
\toprule
\textbf{Variable} & \textbf{Description} \\
\midrule
DATABASE\_URL      & URL de connexion poolée à la base PostgreSQL (Neon). Utilisée par Prisma pour les requêtes applicatives. \\
DIRECT\_URL        & URL de connexion directe (sans pooler). Utilisée uniquement lors des migrations Prisma. \\
SESSION\_SECRET    & Clé secrète (min. 32 caractères) utilisée par Iron Session pour signer et chiffrer les cookies de session. \\
STRIPE\_SECRET\_KEY      & Clé secrète Stripe pour les appels API côté serveur (création de sessions de paiement). \\
STRIPE\_WEBHOOK\_SECRET  & Secret de validation des webhooks Stripe, permettant de vérifier l'authenticité des notifications reçues. \\
RESEND\_API\_KEY   & Clé d'accès au service d'envoi d'e-mails Resend. \\
BLOB\_READ\_WRITE\_TOKEN  & Token d'accès à Vercel Blob pour l'upload et la récupération des partitions PDF. \\
APP\_URL           & URL publique de l'application (ex. \texttt{https://mon-app.vercel.app}). Utilisée pour générer les liens d'activation et de réinitialisation de mot de passe. \\
\bottomrule
\end{tabular}
\caption{Variables d'environnement de l'application}
\end{table}

\newpage

% ─────────────────────────────────────────────
\section{Configuration du pipeline CI/CD}
% ─────────────────────────────────────────────

Le fichier suivant correspond à la configuration GitHub Actions du pipeline
d'intégration continue (\texttt{.github/workflows/ci.yml}).
Il est déclenché à chaque \textit{push} sur la branche \texttt{main} ainsi
qu'à chaque \textit{pull request}, et exécute successivement l'installation
des dépendances et la compilation du projet.

\begin{lstlisting}[language=,caption={.github/workflows/ci.yml}]
name: CI

on:
  push:
    branches: [main]
  pull_request:

jobs:
  build:
    runs-on: ubuntu-latest
    env:
      DATABASE_URL: ${{ secrets.DATABASE_URL }}
      DIRECT_URL: ${{ secrets.DIRECT_URL }}
      STRIPE_SECRET_KEY: ${{ secrets.STRIPE_SECRET_KEY }}

    steps:
      - name: Checkout
        uses: actions/checkout@v4

      - name: Setup pnpm
        uses: pnpm/action-setup@v4

      - name: Setup Node
        uses: actions/setup-node@v4
        with:
          node-version: 20

      - name: Install dependencies
        run: pnpm install --frozen-lockfile

      - name: Build
        run: pnpm build
\end{lstlisting}

\newpage

% ─────────────────────────────────────────────
\section{Configuration des tests unitaires}
% ─────────────────────────────────────────────

Le fichier suivant correspond à la configuration de \textit{Vitest},
l'outil utilisé pour les tests unitaires du projet (\texttt{vitest.config.ts}).

\begin{lstlisting}[language=,caption={vitest.config.ts}]
import { defineConfig } from "vitest/config";
import react from "@vitejs/plugin-react";
import path from "node:path";

export default defineConfig({
  plugins: [react()],
  test: {
    environment: "jsdom",
    globals: true,
    setupFiles: ["./vitest.setup.ts"],
  },
  resolve: {
    alias: {
      "@": path.resolve(__dirname, "."),
    },
  },
});
\end{lstlisting}

Le tableau suivant récapitule les fichiers de tests unitaires rédigés dans le cadre
du projet, organisés par domaine fonctionnel.

\begin{table}[H]
\centering
\renewcommand{\arraystretch}{1.4}
\begin{tabular}{p{8cm} p{6cm}}
\toprule
\textbf{Fichier de test} & \textbf{Fonctionnalité couverte} \\
\midrule
\texttt{test/api/login.test.ts}            & Authentification (identifiants, session) \\
\texttt{test/api/activation.test.ts}       & Activation de compte via lien d'invitation \\
\texttt{test/api/reset-password.test.ts}   & Réinitialisation du mot de passe \\
\texttt{test/api/forgot-password.test.ts}  & Demande de réinitialisation \\
\texttt{test/api/ticket-validate.test.ts}  & Validation des billets (ScanToken) \\
\texttt{test/admin/sessions.test.ts}       & Création et gestion des séances \\
\texttt{test/admin/attendance.test.ts}     & Enregistrement des présences et solde \\
\texttt{test/admin/inscriptions.test.ts}   & Traitement des demandes d'inscription \\
\texttt{test/lib/session.test.ts}          & Gestion des sessions Iron Session \\
\texttt{test/basic.test.ts}               & Tests de base \\
\bottomrule
\end{tabular}
\caption{Récapitulatif des fichiers de tests unitaires (54 tests, 10 fichiers)}
\end{table}

\newpage

% ─────────────────────────────────────────────
\section{Arborescence du projet}
% ─────────────────────────────────────────────

L'arborescence suivante présente la structure générale du projet,
en excluant les répertoires générés automatiquement
(\texttt{node\_modules}, \texttt{.next}, \texttt{generated}, \texttt{coverage}).

\begin{lstlisting}[language=,caption={Structure du projet (71 répertoires, 133 fichiers)},basicstyle=\ttfamily\footnotesize]
.
├── app/
│   ├── activation/          # Activation de compte via lien d'invitation
│   ├── admin/               # Interface d'administration
│   │   ├── courses/         # Gestion des cours
│   │   ├── events/          # Gestion des événements
│   │   ├── inscriptions/    # Traitement des demandes d'inscription
│   │   ├── ressources/      # Gestion des partitions PDF
│   │   ├── sessions/        # Planification des séances
│   │   │   └── [id]/attendance/  # Présences par séance
│   │   ├── students/        # Gestion des élèves
│   │   ├── tarifs/          # Configuration des tarifs
│   │   ├── layout.tsx
│   │   └── page.tsx         # Tableau de bord admin
│   ├── api/                 # API Routes (logique serveur)
│   │   ├── cron/billing/    # Facturation mensuelle automatique
│   │   ├── login/           # Authentification
│   │   ├── resources/[id]/  # Accès protégé aux partitions PDF
│   │   ├── stripe/webhook/  # Webhook de confirmation de paiement
│   │   └── ticket/validate/ # Validation des billets QR code
│   ├── components/          # Composants React partagés
│   ├── contact/             # Formulaire de contact public
│   ├── event/               # Pages publiques des événements
│   │   └── [id]/checkout/   # Flux de paiement Stripe
│   ├── forgot-password/     # Demande de réinitialisation
│   ├── inscription/         # Formulaire d'inscription public
│   ├── login/               # Page de connexion
│   ├── logout/              # Route de déconnexion
│   ├── reset-password/      # Réinitialisation du mot de passe
│   ├── scan/[token]/        # Scan QR code (accès temporaire staff)
│   ├── student/             # Espace personnel de l'élève
│   ├── tarifs/              # Page publique des tarifs
│   └── layout.tsx           # Layout racine de l'application
├── lib/                     # Utilitaires partagés
│   ├── email.ts             # Envoi d'e-mails via Resend
│   ├── prisma.ts            # Client Prisma (singleton)
│   ├── session.ts           # Gestion des sessions Iron Session
│   ├── ticket-pdf.ts        # Génération des billets PDF avec QR code
│   └── translations.ts      # Internationalisation FR/TR
├── prisma/
│   ├── migrations/          # Historique des migrations de base de données
│   └── schema.prisma        # Schéma de la base de données
├── public/                  # Fichiers statiques (images, icônes)
├── scripts/
│   └── seed.ts              # Script d'initialisation de la base de données
├── test/                    # Tests unitaires (Vitest)
│   ├── admin/               # Tests de la logique admin
│   ├── api/                 # Tests des API Routes
│   └── lib/                 # Tests des utilitaires
├── next.config.ts           # Configuration Next.js
├── vitest.config.ts         # Configuration Vitest
└── package.json
\end{lstlisting}

\newpage

% ─────────────────────────────────────────────
\section{Captures d'écran de l'application}
% ─────────────────────────────────────────────

Cette annexe présente les principales interfaces de l'application,
illustrant les vues destinées à l'administrateur et à l'élève.

\begin{figure}[H]
    \centering
    \includegraphics[width=\linewidth]{dashboard-admin.png}
    \caption{Tableau de bord administrateur}
    \label{fig:dashboard-admin}
\end{figure}

\newpage

\begin{figure}[H]
    \centering
    \includegraphics[width=\linewidth]{dashboard-etudiant.png}
    \caption{Espace personnel de l'élève}
    \label{fig:dashboard-etudiant}
\end{figure}

\begin{figure}[H]
    \centering
    \includegraphics[width=\linewidth, height=0.88\textheight, keepaspectratio]{formulaire-inscription.png}
    \caption{Formulaire d'inscription — vue adulte}
    \label{fig:inscription-adulte}
\end{figure}

\newpage

\begin{figure}[H]
    \centering
    \includegraphics[width=\linewidth, height=0.88\textheight, keepaspectratio]{inscription_pour_enfant.png}
    \caption{Formulaire d'inscription — vue parent (inscription d'un mineur)}
    \label{fig:inscription-parent}
\end{figure}

\newpage

% \begin{figure}[H]
%   \centering
%   \includegraphics[width=\linewidth]{presences.png}
%   \caption{Gestion des présences d'une séance}
%   \label{fig:presences}
% \end{figure}

% \begin{figure}[H]
%   \centering
%   \includegraphics[width=\linewidth]{inscriptions.png}
%   \caption{Gestion des demandes d'inscription}
%   \label{fig:inscriptions}
% \end{figure}

\end{document}
